\section{Implementation with Algebraic Effects (OCaml 5.3)}
This section can be read independently; it connects the public API to the
runtime mechanisms and effect handlers.
\tempo{} implements its synchronous model as a handler-driven runtime. The design
keeps the semantics explicit in library code while using OCaml 5.3 algebraic
effects to capture and schedule continuations. This section focuses on how the
effects-based implementation realizes the \tempo{} API.

\paragraph{Implementation map}
For readers skimming the code, the implementation can be read in three passes:
(1) the effect declarations in \texttt{lib/tempo\_types.ml}; (2) the scheduler
loop and task queues in \texttt{lib/tempo.ml}; (3) the handlers for each effect
in the same file. The most important invariant to keep in mind is that all
continuations are resumed by enqueueing tasks, never by direct calls, which
preserves the instant discipline.

\paragraph{Scheduler state and task queues}
The runtime maintains three task collections: a FIFO queue of tasks runnable
in the current instant, a list of tasks scheduled for the next instant, and a
list of tasks blocked on unsatisfied guards. Each task carries its guards and
preemption tokens, and tasks are grouped into logical threads so that
\texttt{join} can synchronize on thread completion. A registry of signals
tracks presence, current value, and waiting continuations.

The scheduler state also stores debug counters (signal ids, task ids, step and
instant counters) and a per-thread bookkeeping table. Each thread state keeps
an active task count and a list of waiters to resume when the thread completes.
This enables \texttt{join} to synchronize without global locks or OS threads.

The scheduler is single-threaded and deterministic: all decisions are made in
one runtime loop. This makes it easier to reason about the ordering of effects
within an instant, but it also means that the implementation must explicitly
track every deferred continuation to avoid hidden control flow.

Tasks move between queues as follows: when a task performs an effect that
defers it (e.g., \texttt{await} or \texttt{pause}), its continuation is wrapped
as a new task and enqueued in \texttt{next\_instant}. When a task is blocked by
guards, it is placed in the \texttt{blocked} list and also registered in the
missing signals' \texttt{guard\_waiters}. At instant end, all blocked tasks are
moved to \texttt{next\_instant}, preserving the rule that absent guards only
become observable in the next instant.

The scheduler also maintains a clear quiescence condition: an instant ends
only when the current queue is empty. This ensures that all reactions to
present signals have been taken before finalization, aligning with the
delayed-absence semantics.

\paragraph{Task lifecycle and bookkeeping}
Creating a task registers it under a thread id and increments the active count.
When a task finishes, the active count is decremented; if it reaches zero, any
waiters are resumed. Tasks carry a \emph{queued} flag to avoid double enqueue,
a \emph{blocked} flag to control guard rescheduling, and a list of kill tokens
that determine whether the task is still alive. This lifecycle is essential to
the deterministic semantics because it ensures tasks are resumed only at
well-defined instant boundaries.

\paragraph{Threads, \texttt{fork}, and \texttt{join}}
A \texttt{fork} creates a new logical thread with its own task counter. The
parent task continues immediately, while the child task is scheduled in the
current instant. \texttt{join} captures the caller's continuation and registers
it in the target thread's waiter list. When the thread's active counter reaches
zero, all join waiters are resumed in the next instant. This design allows
structured concurrency without OS threads, while still making completion
observable at instant boundaries.

\begin{lstlisting}[style=tempo]
(* Pseudo-code for join handling *)
on Join thread_id:
  if thread completed then
    enqueue_next (continue k ())
  else
    add waiter (fun () -> enqueue_next (continue k ()))
\end{lstlisting}

\paragraph{Invariants}
The scheduler maintains several invariants that are critical for correct
behavior:
\begin{itemize}
\item \textbf{No direct resumption}: every continuation becomes a task and is
  enqueued through the scheduler queues.
\item \textbf{No re-entrancy}: a task is never scheduled while it is already
  running (\texttt{queued} flag).
\item \textbf{Guard discipline}: a blocked task appears in guard\_waiters of
  every missing guard and nowhere else.
\item \textbf{Kill discipline}: a task is scheduled only if all its kill tokens
  are alive.
\end{itemize}
Most of the implementation complexity exists to preserve these invariants in
the presence of nested guards, nested preemption, and parallel composition.

\paragraph{Failure modes and debug hooks}
Several runtime checks enforce the synchronous discipline:
\begin{itemize}
\item \textbf{Multiple emits on event signals}: a second emission in the same
  instant raises an exception because the event semantics require a single
  value per instant.
\item \textbf{Present-without-value}: if a signal is marked present but has no
  stored value, the runtime fails fast to avoid undefined behavior.
\item \textbf{Dead kill tokens}: tasks protected by aborted kill tokens are
  discarded silently; this is intentional and implements weak preemption.
\end{itemize}
Debug counters (task ids, signal ids, instant/step counters) and optional logs
make it possible to correlate observable outputs with internal queue states.
These hooks are crucial when diagnosing unexpected blocking or when confirming
that a program respects the instant discipline.

\paragraph{Core runtime types}
The implementation centers on a small set of runtime records. The following
summary aligns with the concrete types in \texttt{lib/tempo\_types.ml} and
\texttt{lib/tempo.ml}.
\begin{table}[h]
\centering
\begin{tabularx}{\linewidth}{lY}
\hline
Type & Role and fields (informal) \\ \hline
\texttt{task} & \texttt{t\_id}, \texttt{thread}, \texttt{guards}, \texttt{kills},
\texttt{run}, plus \texttt{queued}/\texttt{blocked} flags. \\
\texttt{signal\_core} & \texttt{s\_id}, \texttt{present}, \texttt{value},
\texttt{awaiters}, \texttt{guard\_waiters}, and \texttt{kind} (event or
aggregate with \texttt{combine}/\texttt{initial}). \\
\texttt{thread\_state} & \texttt{active}, \texttt{completed}, \texttt{waiters}. \\
\texttt{scheduler\_state} & \texttt{current}, \texttt{next\_instant},
\texttt{blocked}, \texttt{signals}, \texttt{thread\_counter}, \texttt{threads},
and \texttt{debug} counters. \\ \hline
\end{tabularx}
\caption{Key runtime structures used by the scheduler.}
\end{table}

\paragraph{Declared effects}
The reactive operations are exposed as algebraic effects in
\texttt{lib/tempo\_types.ml}. The code below mirrors the actual declaration:
\begin{lstlisting}[style=tempo]
type _ Effect.t +=
  | New_signal : unit -> ('a, 'a, event) signal_core Effect.t
  | New_signal_agg :
      'agg * ('agg -> 'emit -> 'agg)
      -> ('emit, 'agg, aggregate) signal_core Effect.t
  | Emit : ('emit, 'agg, 'mode) signal_core * 'emit -> unit Effect.t
  | Await : ('emit, 'agg, 'mode) signal_core -> 'agg Effect.t
  | Await_immediate : ('a, 'a, event) signal_core -> 'a Effect.t
  | Pause : unit Effect.t
  | Fork : (unit -> unit) -> thread Effect.t
  | Join : thread -> unit Effect.t
  | With_guard :
      ('emit, 'agg, 'mode) signal_core * (unit -> unit) -> unit Effect.t
  | With_kill : kill * (unit -> unit) -> unit Effect.t
\end{lstlisting}
Each effect corresponds to a single scheduler action: \texttt{New\_signal}
allocates a fresh signal; \texttt{Emit} publishes a value and wakes waiters;
\texttt{Await}/\texttt{Await\_immediate} suspend and resume continuations;
\texttt{Pause} shifts execution to the next instant; \texttt{Fork}/\texttt{Join}
manage logical threads; and \texttt{With\_guard}/\texttt{With\_kill} implement
guarding and weak preemption.

\paragraph{From effect to scheduler (worked example)}
Consider \texttt{Await} as implemented in \texttt{lib/tempo.ml}. When a task
performs \texttt{Await s}, the handler captures the continuation and registers
it in \texttt{s.awaiters}. If \texttt{s} is already present and is an event
signal, the continuation is scheduled for the next instant; otherwise it is
stored as a waiter and resumed when \texttt{s} becomes present. The critical
behavior is that resumption is \emph{always} scheduled through the scheduler
queue, never executed directly in the current stack, which preserves the
instant boundary discipline.

In the excerpt below, the helper \texttt{mk\_task} wraps a continuation into a
runnable task. \texttt{enqueue\_next} schedules it for the next instant, and
\texttt{continue k v} resumes the captured continuation with value \texttt{v}.

\begin{lstlisting}[style=tempo]
(* lib/tempo.ml: effect handling for Await *)
| effect (Await s), k ->
    let enqueue_resume v =
      let t' = mk_task st t.thread t.guards t.kills
        (fun () -> continue k v) in
      enqueue_next st t'
    in
    let register_waiter () =
      let resume v =
        let t' = mk_task st t.thread t.guards t.kills
          (fun () -> continue k v) in
        enqueue_next st t'
      in
      s.awaiters <- resume :: s.awaiters
    in
    if s.present then
      match s.kind with
      | Event_signal ->
          (match s.value with
           | Some v -> enqueue_resume v
           | None -> failwith "present but no value")
      | Aggregate_signal _ ->
          register_waiter ()
    else
      register_waiter ()
\end{lstlisting}

In terms of observable behavior, this means that an \texttt{await} never
resumes in the current instant; the continuation is always deferred to the
next instant, ensuring that presence and absence are observed only at instant
boundaries. This pattern is representative of how effects encode the scheduler
policy in user space.

\begin{table}[h]
\centering
\begin{tabularx}{\linewidth}{lY}
\hline
Effect & Scheduler action and instant semantics \\ \hline
\texttt{New\_signal} & Allocate signal state in the current instant. \\
\texttt{Emit} & Mark present and wake awaiters (possibly now). \\
\texttt{Await} & Register continuation; resume next instant. \\
\texttt{Await\_immediate} & Resume now if present; otherwise register waiter. \\
\texttt{Pause} & Defer continuation to the next instant. \\
\texttt{Fork} & Spawn child task in current instant. \\
\texttt{Join} & Suspend until thread completes, then resume. \\
\texttt{With\_guard} & Attach guard set; block if guards absent. \\
\texttt{With\_kill} & Register cleanup; abort at instant end on kill. \\ \hline
\end{tabularx}
\caption{Effect-to-scheduler mapping.}
\end{table}

\paragraph{Instant loop and signal finalization}
Execution proceeds by instants. During an instant, the scheduler repeatedly
dequeues a runnable task, checks whether its guards are satisfied, and either
runs it or moves it to the blocked list. At the end of the instant, the
runtime finalizes all signals: event and aggregate signals are reset to
absent, guard waiters are cleared, and aggregate-signal awaiters are resumed
with the accumulated value. Blocked tasks are moved to the next-instant list,
and the next instant begins with the tasks that remain alive.

More precisely, each step records timing metrics, logs queue snapshots, and
updates counters for reproducibility. When the runnable queue becomes empty,
the instant ends. All blocked tasks are re-queued for the next instant, which
implements the delayed-absence discipline: a guard that was absent remains
absent for the entire instant, and the task can only retry in the next one.

Re-queuing does not lose guard information: the blocked task still carries its
guard list, so when it is dequeued in the next instant it is checked again. If
some guards are still absent, it is re-registered in the corresponding
\texttt{guard\_waiters} lists. This is why clearing \texttt{guard\_waiters} at
instant end is safe: the lists are reconstructed as tasks are revisited.

\paragraph{Queue discipline and fairness}
The current queue is FIFO, which yields a simple notion of fairness within an
instant: tasks are resumed in the order they are enqueued. This ordering is an
implementation choice rather than a semantic guarantee, and observable results
should not depend on it. The ordering is nevertheless useful for debugging and
reproducible traces, and it aligns with the cooperative scheduling model of
FairThreads.

To avoid starvation across instants, any task that remains blocked is
reconsidered in the next instant. This provides a coarse-grained fairness
property: if a guard becomes present in some future instant, the corresponding
task will be eligible to run at the start of that instant.

\paragraph{Instant loop pseudo-code}
\begin{lstlisting}[style=tempo]
while current_queue not empty do
  t <- pop current_queue
  if not kills_alive t then discard t
  else if guard_ok t.guards then run t
  else
    block t on missing_guards
    register t in guard_waiters of each missing guard
end
finalize_signals ()
move blocked tasks to next_instant
swap current_queue and next_instant
\end{lstlisting}
This sketch omits bookkeeping details (thread counters, debug logs) but
captures the control structure of the scheduler.

The ordering above is deliberate. Kill checks precede guard checks so that a
preempted task is never allowed to run or block again. Guard checks precede
execution to preserve the semantics that guards are blocking conditions, not
filters applied after a step. If these checks were reversed, a task might run
in an instant where its guard is absent or continue after being preempted,
which would violate the intended weak-preemption semantics.

\begin{figure}[t]
\centering
\resizebox{0.80\columnwidth}{!}{%
\footnotesize
\begin{tikzpicture}[
  node distance=10mm,
  box/.style=tempo,
  draw,
  rounded corners,
  align=center,
  minimum width=30mm,
  minimum height=8mm,
  decision/.style=tempo,
  draw,
  align=center,
  aspect=2,
  inner sep=1pt,
  arrow/.style=tempo
]
\node[box] (start) {Start instant};
\node[box, below=of start] (dequeue) {Dequeue task};
\node[decision, below=of dequeue] (guards) {Guards\\satisfied?};
\node[box, left=of guards] (block) {Block task\\on guards};
\node[box, right=of guards] (handle) {Handle task\\effects};
\node[decision, below=of guards] (empty) {Queue\\empty?};
\node[box, below=of empty] (finalize) {Finalize signals\\and move blocked};
\node[box, below=of finalize] (next) {Next instant};

\draw[arrow] (start) -- (dequeue);
\draw[arrow] (dequeue) -- (guards);
\draw[arrow] (guards) -- node[above]{no} (block);
\draw[arrow] (guards) -- node[above]{yes} (handle);
\draw[arrow] (block) |- (empty);
\draw[arrow] (handle) |- (empty);
\draw[arrow] (empty) -- node[right]{yes} (finalize);
\draw[arrow] (empty) -- ++(20mm,0mm) node[right]{no} |- (dequeue);
\draw[arrow] (finalize) -- (next);
\end{tikzpicture}%
}
\caption{High-level flowchart of the instant-based scheduler.}
\end{figure}

\paragraph{Execution trace (guard + watch)}
Consider the running example from Section~\ref{ex:guard-watch}. In instant $t$,
one task is guarded by signal \texttt{g}, and one is wrapped in \texttt{watch s}.
At the start of the instant, both tasks are in the current queue. The guarded
task is blocked because \texttt{g} is absent and is registered in
\texttt{g.guard\_waiters}. The watched task runs; if it emits \texttt{s} during
$t$, the guardian aborts the kill token and the watched task is prevented from
continuing into $t{+}1$. At instant end, guard waiters are cleared, blocked
tasks are moved to the next-instant list, and the system advances. This trace
illustrates how guard blocking and weak preemption interact with the instant
boundary discipline.

\begin{lstlisting}[style=tempo]
let g = new_signal () in
let s = new_signal () in
let guarded = (fun () -> when_ g (fun () -> emit s ())) in
let watched = (fun () -> watch s (fun () -> pause (); emit s ())) in
parallel [guarded; watched]
\end{lstlisting}
The trace tables below show the corresponding queue, signal, and observable
states across instants.

\begin{table}[h]
\centering
\begin{tabularx}{\linewidth}{lY}
\hline
Instant & Queue/blocked/signals snapshot \\ \hline
$t$ start & Current: \texttt{guarded}, \texttt{watched}; blocked: $\varnothing$;
signals: $g=\bot$, $s=\bot$. \\
$t$ mid & \texttt{guarded} blocked on $g$; \texttt{watched} runs and emits $s=()$;
kill token marked. \\
$t$ end & Guard waiters cleared; blocked $o$ next instant; $s$ observed
present at boundary, then reset; \texttt{watched} will not continue into
$t{+}1$. \\ \hline
\end{tabularx}
\caption{Concrete execution trace for guard + watch (example above).}
\end{table}

\begin{table}[h]
\centering
\begin{tabularx}{\linewidth}{lY}
\hline
Instant & Observable signals and outputs \\ \hline
$t$ & $s=()$ becomes present due to \texttt{watched}; no output emitted. \\
$t{+}1$ & \texttt{guarded} remains blocked if $g=\bot$; no output emitted. \\
\hline
\end{tabularx}
\caption{Two-instant observable trace for the example program.}
\end{table}

\paragraph{Signal operations}
Signal creation effects allocate runtime state and register signals in the
global registry. \texttt{Emit} marks a signal present and stores its value.
For event signals, multiple emissions in the same instant are rejected, while
aggregate signals fold successive emissions into an accumulator. Awaiting a
signal registers a continuation: \texttt{await} always resumes at the next
instant, while \texttt{await\_immediate} can resume in the current instant when
the event signal is already present. This distinction matches the delayed
absence discipline while still allowing immediate reactions when safe.
Figure~\ref{fig:signal-state} summarizes the signal state evolution within a
single instant.

Event signals enforce a single-emission rule. The implementation checks
\texttt{s.present} and raises an exception on a second emission in the same
instant. This runtime check enforces the synchronous discipline in the absence
of static typing rules.

Signal waiters are stored as callbacks (continuations wrapped as functions).
For event signals, each emission registers resumption through the scheduler;
\texttt{await} resumes in the next instant, while \texttt{await\_immediate} may
resume in the current instant if the signal is already present. For aggregate
signals, emissions only update the accumulator; awaiters are resumed once at
instant end with the aggregated value. This makes the point of observation
explicit and keeps the aggregation order-insensitive.

Aggregate signals are delivered only once per instant. Emissions update the
accumulator, but awaiters are resumed at instant end with the accumulated
value. This preserves determinism because the order of emissions within the
instant does not affect the observable schedule, only the aggregated result.
If the combine function is not associative/commutative, the aggregate value may
still depend on intra-instant evaluation order, which is intentionally left
unspecified; this is why the API recommends pure, order-insensitive combines.

The finalization step can be summarized as follows:
\begin{lstlisting}[style=tempo]
for each signal s in registry do
  if s is aggregate then
    resume all awaiters with s.value
  clear s.awaiters and s.guard_waiters
  reset s.present and s.value
\end{lstlisting}
In the actual implementation, event-signal awaiters have already been resumed
on emission, whereas aggregate awaiters are resumed in this finalization phase.

\begin{figure}[t]
\centering
\resizebox{0.92\columnwidth}{!}{%
\begin{tikzpicture}[
  node distance=10mm,
  box/.style=tempo,
  draw,
  rounded corners,
  align=center,
  minimum width=36mm,
  minimum height=8mm,
  arrow/.style=tempo
]
\node[box] (start) {Start instant:\\present=false, value=None};
\node[box, below=of start] (emit) {Emissions:\\present=true, value updated};
\node[box, below=of emit] (awaiters) {Awaiters:\\event resume now\\aggregate waits};
\node[box, below=of awaiters] (finalize) {End instant:\\deliver aggregate,\\clear guard\_waiters};
\node[box, below=of finalize] (reset) {Reset state:\\present=false, value=None};

\draw[arrow] (start) -- (emit);
\draw[arrow] (emit) -- (awaiters);
\draw[arrow] (awaiters) -- (finalize);
\draw[arrow] (finalize) -- (reset);
\end{tikzpicture}%
}
\caption{Signal state evolution within an instant.}
\label{fig:signal-state}
\end{figure}

\paragraph{Guards and blocking}
Guarded execution is implemented by attaching a list of guard signals to each
task. If a guard is missing, the task is blocked and registered in the guard
waiter lists of the missing signals. When a signal is emitted, its guard
waiters are rechecked: tasks whose guards are now all present are moved back to
the runnable queue, otherwise they remain blocked on the remaining missing
guards. This implements the disjunctive behavior of nested guards while
preserving determinism.

Guard checking is centralized: \texttt{guard\_ok} verifies all guards are
present, while \texttt{missing\_guards} returns the absent subset. This keeps
guard logic explicit and ensures the scheduler, rather than user code, controls
when a blocked task can resume.

To avoid repeated scans, each signal maintains its own \texttt{guard\_waiters}
list. When a signal becomes present, only tasks that were blocked on that
signal are reconsidered. Tasks may remain blocked if other guards are still
absent, and they are re-registered on the remaining missing guards. This local
indexing keeps guard rescheduling proportional to the number of blocked tasks
that are directly affected by the emission.

\paragraph{Guard disjunction by nesting}
Nested \texttt{when\_} expressions behave as a conjunction at the level of
guard satisfaction: a task blocked on multiple guards remains blocked until
\emph{all} guards are present in the same instant. If a program wants a logical
disjunction (guard $g_1$ \emph{or} $g_2$), it must use parallel guarded
branches, which creates distinct tasks each blocked on a single guard. This
mechanism follows directly from the guard list representation and the
centralized \texttt{guard\_ok} predicate.

\paragraph{Effect \texttt{With\_guard} and \texttt{when\_}}
The \texttt{when\_} combinator is implemented using the \texttt{With\_guard}
effect. Performing \texttt{With\_guard (s, body)} creates a new task whose
guard list is extended with \texttt{s}. If \texttt{s} is absent, the task is
blocked and registered in \texttt{s.guard\_waiters}. Once \texttt{s} becomes
present, the scheduler rechecks the guard list and enqueues the task in the
current instant. This means that \texttt{when\_} does not execute \texttt{body}
until the guard is satisfied, and the scheduling decision remains centralized.
In the excerpt below, \texttt{Any s} wraps a typed signal into an existential
for the guard list, and \texttt{enqueue\_now} schedules the guarded task in the
current instant.
\begin{lstlisting}[style=tempo]
(* lib/tempo.ml: effect handling for With_guard *)
| effect (With_guard (s, body)), k ->
    let t' =
      mk_task st t.thread (Any s :: t.guards) t.kills
        (fun () ->
          body ();
          enqueue_now st (mk_task st t.thread t.guards t.kills
            (fun () -> continue k ())))
    in
    enqueue_now st t'
\end{lstlisting}

\paragraph{Guard resumption trace}
Consider a task guarded by $g_1$ and $g_2$. At instant $t$, $g_1$ is present
but $g_2$ is absent. The task is blocked and registered in $g_2.guard\_waiters$
only; it is not re-checked until $g_2$ becomes present. When $g_2$ is emitted
in instant $t{+}1$, the scheduler revisits the task, finds all guards present,
and enqueues it in the current queue. This preserves the intended "all guards
in the same instant" semantics.

\paragraph{Effect \texttt{With\_kill} and \texttt{watch}}
Weak preemption relies on the \texttt{With\_kill} effect. A kill token stores a
cleanup continuation that is invoked when the token is aborted. The \texttt{watch}
combinator creates a shared kill token for the protected body and a guardian
task that waits on the watched signal. When the signal is emitted, the guardian
aborts the shared token, which triggers the cleanup and prevents further
resumption of the protected task.

Concretely, \texttt{watch} runs the body under \texttt{With\_kill} and installs
a cleanup continuation that schedules the parent continuation at the appropriate
instant boundary. The scheduler checks \texttt{kills\_alive} before enqueuing a
task, so any task protected by an aborted kill token is discarded. This yields
weak preemption: the body can finish the current instant, but it will not
continue in subsequent instants once the watched signal is present.
In the excerpt below, \texttt{continue\_now} resumes the parent continuation in
the current instant, while \texttt{continue\_later} defers it to the next
instant. The cleanup stored in \texttt{kk.cleanup} chooses the safe resumption
point once the kill token is aborted.
\begin{lstlisting}[style=tempo]
(* lib/tempo.ml: effect handling for With_kill *)
| effect (With_kill (kk, body)), k ->
    let continue_now () =
      kk.cleanup <- None;
      enqueue_now st (mk_task st t.thread t.guards t.kills
        (fun () -> continue k ()))
    and continue_later () =
      kk.cleanup <- None;
      enqueue_next st (mk_task st t.thread t.guards t.kills
        (fun () -> continue k ()))
    in
    kk.cleanup <- Some continue_later;
    let runner () =
      body ();
      continue_now ()
    in
    let t' = mk_task st t.thread t.guards (kk :: t.kills) runner in
    enqueue_now st t'
\end{lstlisting}

\paragraph{Worked walkthrough: \texttt{watch}}
Consider a body \texttt{B} under \texttt{watch s}. At instant $t$, the guardian
awaits \texttt{s} while \texttt{B} runs. If \texttt{s} is emitted during $t$,
the guardian aborts the shared kill token and installs cleanup to prevent
\texttt{B} from continuing at $t{+}1$. If \texttt{s} is not emitted, \texttt{B}
continues normally in the next instant. This aligns with weak preemption: the
body can complete the current instant but is cut off before the next one.

\paragraph{Weak preemption with kill tokens}
Weak preemption is modeled through kill tokens. A \texttt{watch} spawns a
guardian that waits on the watched signal and aborts the protected computation
at the end of the instant in which the signal is emitted. Internally, the
runtime stores cleanup continuations for kill tokens so that aborting a token
can schedule the appropriate continuation either immediately or in the next
instant. This approach encodes preemption without violating the single-instant
deterministic semantics.

Kill tokens also ensure that tasks are not resumed after preemption. The
scheduler checks that all kill tokens are alive before enqueuing a task; if any
token is dead, the task is discarded and the thread bookkeeping is updated.
Because preemption is weak, application code sometimes adds an explicit guard
(e.g., a \texttt{handled} signal) to prevent simultaneous emissions in the same
instant; this pattern is illustrated in the usage examples.

\paragraph{Kill tokens and nested \texttt{watch}}
Nested \texttt{watch} constructs install multiple kill tokens. The scheduler
checks that \emph{all} kill tokens are alive before running a task. This means
any inner \texttt{watch} that aborts its token is sufficient to prevent the
task from continuing, which realizes the intended disjunction of stopping
conditions. Kill tokens do not cancel the current instant; they only prevent
the continuation from being enqueued after the instant boundary, which
preserves weak preemption.

\paragraph{Kill cleanup timing}
The cleanup stored in a kill token is intentionally deferred with respect to
observable behavior: if the watched signal is emitted in instant $t$, the
protected computation can still finish its current instant but will not be
resumed in $t{+}1$. Internally, the handler may schedule the cleanup either
immediately or for the next instant, but the user-visible effect is always
weak preemption at the instant boundary.

\paragraph{Host I/O hooks}
The entry point \texttt{execute} creates dedicated input and output signals.
Before each step, host input can emit on the input signal; after the step, the
output signal is checked and delivered to the host. This keeps external I/O
separate from synchronous reactions while preserving the instant structure.

Internally, \texttt{execute} installs a handler for all effects and runs a
driver loop that alternates between (1) injecting host input for the instant,
(2) running the scheduler until quiescence, and (3) collecting the output
signal. This design keeps the host side free of continuation manipulation: the
host only sees a stream of inputs and outputs at instant boundaries, which is
consistent with the synchronous model.

\begin{lstlisting}[style=tempo]
loop:
  call input_hook to emit on input signal
  run scheduler until current queue is empty
  read output signal at instant end
\end{lstlisting}
Outputs are therefore observed only at instant boundaries, which matches the
delayed-absence discipline and prevents mid-instant I/O from leaking
implementation order.

For readability, helper functions used in the excerpts are summarized below
(see Section~\ref{sec:helper-glossary}).

\paragraph{Runtime helper glossary}\label{sec:helper-glossary}
\begin{itemize}
\item \texttt{mk\_task}: wrap a continuation into a runnable task record.
\item \texttt{enqueue\_now}: schedule a task in the current instant queue.
\item \texttt{enqueue\_next}: schedule a task for the next instant.
\item \texttt{continue}: resume a captured continuation with a value.
\item \texttt{Any}: existential wrapper to store signals in guard lists.
\item \texttt{kills\_alive}: check that all kill tokens of a task are alive.
\end{itemize}

\subsection{Algebraic Effects in Detail}
	\tempo{} encodes the reactive primitives as algebraic effects and interprets them
with a deep handler. Each effect represents a semantic action (e.g., emit,
await, pause), and the handler decides whether to resume the continuation now
or to re-schedule it for a later instant. This design makes the control flow
explicit and keeps the scheduler in user space rather than in a custom runtime.

Effect capture preserves the current continuation, which is then wrapped into a
task structure. For example, \texttt{await} captures the continuation, stores a
resume function in the signal's waiter list, and ensures that the continuation
is resumed at the beginning of the next instant. \texttt{await\_immediate}
selects between immediate resumption (if the signal is already present) and
registration as a waiter (if it is not). Similarly, \texttt{pause} captures the
continuation and places it in the next-instant list.

Signal creation effects allocate fresh signal records and return them to the
continuation. \texttt{Emit} updates signal state and resumes any registered
awaiters, potentially in the current instant. Guard effects update the task's
guard set and may block the task if the guards are not satisfied. Preemption
effects (\texttt{With\_kill}) store cleanup continuations that are invoked when
the kill token is aborted, ensuring that termination is coordinated with the
instant boundary.

This effect-based organization provides a clear separation between the
operational model (captured by the handler) and the user-facing API (exposed as
effectful functions). It also enables experimentation with scheduling
strategies without changing the OCaml runtime or the language semantics.

\paragraph{Handler pseudo-code}
\begin{lstlisting}[style=tempo]
handle task =
  match task() with
  | effect (Emit (s, v)) k ->
      mark_present s v;
      resume_awaiters s;
      enqueue_now (continue k ())
  | effect (Await s) k ->
      register_waiter s (fun v -> enqueue_next (continue k v))
  | effect (Await_immediate s) k ->
      if present s then enqueue_now (continue k (value s))
      else register_waiter s (fun v -> enqueue_now (continue k v))
  | effect Pause k ->
      enqueue_next (continue k ())
  | effect (With_guard (s, body)) k ->
      update_guards s;
      enqueue_now (fun () -> body (); continue k ())
  | effect (With_kill (kk, body)) k ->
      install_cleanup kk (fun () -> enqueue_next (continue k ()));
      enqueue_now (fun () -> body (); continue k ())
  | effect (Fork f) k ->
      enqueue_now f; enqueue_now (continue k ())
  | effect (Join t) k ->
      add_join_waiter t (fun () -> enqueue_now (continue k ()))
\end{lstlisting}

\subsection{Detailed Effect Traces}
This subsection provides step-by-step traces for key primitives. The goal is
to make explicit where continuations are stored and how they move across
queues.

\paragraph{Trace: \texttt{await\_immediate}}
Suppose \texttt{s} is an event signal. If \texttt{s} is already present, the
continuation is resumed in the current instant via \texttt{enqueue\_now};
otherwise it is registered in \texttt{s.awaiters} for resumption at the next
instant. The behavior is summarized below.
\begin{table}[h]
\centering
\begin{tabularx}{\linewidth}{lY}
\hline
Case & Scheduler action \\ \hline
\texttt{s.present=true} & Enqueue continuation now; current instant resumes. \\
\texttt{s.present=false} & Store waiter; resume at next instant. \\ \hline
\end{tabularx}
\caption{Scheduler actions for \texttt{await\_immediate}.}
\end{table}

\paragraph{Trace: \texttt{pause}}
\texttt{pause} captures the current continuation and enqueues it for the next
instant. The task does not run again in the current instant, which is the
primary mechanism for separating instants in user code.

\paragraph{Trace: \texttt{parallel}}
\texttt{parallel} spawns multiple tasks in the current instant and waits for
their completion. It relies on \texttt{fork} under the hood: each branch is a
child task, and the parent waits until the child thread's active count reaches
zero. The effect-level view is therefore:
\begin{enumerate}
\item Spawn a child task for each branch.
\item Continue parent task to a \texttt{join}-like waiting point.
\item Resume parent once all branches terminate.
\end{enumerate}

\paragraph{Trace: \texttt{watch}}
The following sketch shows the lifecycle of a watched task.
\begin{table}[h]
\centering
\begin{tabularx}{\linewidth}{lY}
\hline
Step & Scheduler action \\ \hline
Start $t$ & Spawn guardian awaiting $s$; run body task. \\
$s$ emitted & Guardian aborts kill token; installs cleanup. \\
End $t$ & Body may complete; continuation not enqueued for $t{+}1$. \\
$t{+}1$ & Body is discarded if kill token is dead. \\ \hline
\end{tabularx}
\caption{Execution trace for \texttt{watch}.}
\end{table}

\paragraph{Trace: \texttt{when\_}}
The guard handler creates a task with an extended guard list. If the guard is
absent, the task is blocked; if the guard is present, the task is enqueued
immediately. The observable result is that \texttt{body} only begins in an
instant where the guard is present.

\subsection{Code Map: from API to Runtime}
The table below provides a roadmap from surface-level functions to the runtime
mechanisms that implement them. This can be used as a guide when reading the
source.
\begin{table}[h]
\centering
\begin{tabularx}{\linewidth}{lY}
\hline
API function & Runtime mechanism \\ \hline
\texttt{new\_signal} & \texttt{New\_signal} effect, allocates signal\_core. \\
\texttt{new\_signal\_agg} & \texttt{New\_signal\_agg} effect, sets combine/init. \\
\texttt{emit} & \texttt{Emit} effect, updates signal and wakes waiters. \\
\texttt{await} & \texttt{Await} effect, registers waiter for next instant. \\
\texttt{await\_immediate} & \texttt{Await\_immediate} effect, immediate or next. \\
\texttt{pause} & \texttt{Pause} effect, enqueue for next instant. \\
\texttt{when\_} & \texttt{With\_guard} effect, extend guard list. \\
\texttt{watch} & \texttt{With\_kill} effect, install kill token/cleanup. \\
\texttt{parallel} & \texttt{fork}/\texttt{join}, thread counters. \\
\texttt{execute} & Scheduler loop + host I/O hooks. \\ \hline
\end{tabularx}
\caption{Mapping from API to runtime mechanisms.}
\end{table}
